\section{Introduction}

A constitutional polity has two halves of authority. The first half is
the admission predicate that decides whether an incoming receipt belongs
to the polity's history; the second half is the positive enforcement act
the polity issues when its constitution requires response to an admitted
fact. The amendment lifecycle has been formalized on the first half: a
candidate constitution becomes the active constitution only when the
amendment witness carries a Lean term establishing that every receipt
admitted by the new predicates was also admitted by the old. The
response side has been left to the admission predicate, which means an
enforcement act has had no positive type-level discipline of its own:
the only thing the type system enforced was that the act admitted, not
that the act terminated, reverted, or carried a counterparty's signature.

The reversible-action discipline extends the type-level rule from
amendment admission to enforcement response, for the class of actions
whose effects are reversible by an inverse executor. A reversible
executive action is a type whose constructor refuses to build a witness
without a positive time-to-live duration and an action class tag; the
inverse executor produces a rollback receipt that closes the action
chain by referencing the original execution receipt and the executor key
that signed the rollback. A destructive action, defined by the absence
of an inverse executor, reduces admission to bilateral predicate
intersection between a device polity that hosts the action and an
operator polity that authorizes it; the reduction specializes the prior
treaty intersection theorem \cite{programmableSovereigntyAnonymous} to
the destructive class.

The load-bearing composition claim relates the response side to the
amendment side. A TTL-bounded amendment is a constitutional delta paired
with a positive TTL duration: the new constitution is active inside the
window, and the old constitution returns when the window expires. Such
amendments form a chain when issued sequentially. Each chain step
carries the same per-step backward-refinement witness the amendment
lifecycle already requires. The composition theorem states that for any
instant during the chain's lifetime, the active constitution refines the
baseline at that instant, regardless of whether the instant falls inside
or outside any given amendment's TTL window. The auto-reversion at TTL
expiry is exactly the structural mechanism that keeps the chain's
trajectory pinned to the baseline rather than drifting predicate-by-
predicate into a ratchet.

The contributions are four artifacts that follow the same falsifiability
discipline as the prior substrate
\cite{programmableSovereigntyAnonymous}.

\begin{itemize}
  \item A formal model relating the executive-action type, its TTL
    invariant, its optional rollback receipt, and the TTL-bounded
    amendment chain (Section~\ref{sec:model-ra}).
  \item A load-bearing composition theorem
    \thm{ttl_bounded_amendment_chain_preserves_baseline} and a
    supporting rollback-receipt admission theorem
    \thm{rollback_admission_composes_with_refinement}, both mechanized
    in a paper-local Lean model on an explicit finite admission domain,
    with two definitional bridges retained to anchor the type-level
    discipline (Section~\ref{sec:model-ra}).
  \item A proposed endpoint substrate instance in which four reversible
    action variants would have inverse executors with content-hash-gated
    rollback constructors and a four-receipt grammar (request, execution,
    rollback, acknowledgement) records each step of the action chain
    (Section~\ref{sec:substrate-ra}).
  \item An evaluation plan that names the reversible-path harnesses,
    fixture corpora, destructive class, missing TTL scheduler, and stub
    endpoint sensor as gaps rather than reporting unshipped measurements
    (Section~\ref{sec:evaluation-ra}).
\end{itemize}

The sovereignty claim is bounded. Sovereignty over reversible action
covers the polity's receipt log, not the physical world; a rollback
receipt is a statement about that log referencing an inverse executor's
signed assertion that the OS operation completed. The receipt-log
discipline is the falsifiable part; the physical reversibility is an
empirical claim about specific executor pairs (rename and inverse
rename, signal stop and signal continue) that the receipt schema
records without itself proving.
